% TOP_PROBLEM: {"id": 5300044, "title": "Invariant line fields on Julia sets", "category_id": 11, "category": {"id": 11, "name": "dynamical_systems", "display_name": "Dynamical Systems", "description": "Problems about long-term behavior of deterministic systems, Hamiltonian dynamics, and periodic orbits.", "slug": "dynamical-systems", "order_index": 11, "created_at": "2026-07-31T15:26:25.671Z"}, "status": "open", "problem_number": "AMR-052-0044", "set_id": 13}
% FIRST_POSTED: 2026-10-01
% ENTRY_KIND: statement_only
% UnsolvedMath ID 5300044; source catalogue code AMR-052-0044
% !TeX program = lualatex
% Definition and sources only; this file is not an LLM solution attempt.
\documentclass[11pt,a4paper]{article}
\usepackage[margin=25mm]{geometry}
\usepackage{fontspec}
\setmainfont{lmroman10-regular.otf}[BoldFont=lmroman10-bold.otf,
 ItalicFont=lmroman10-italic.otf,BoldItalicFont=lmroman10-bolditalic.otf]
\usepackage{amsmath,amssymb,mathtools}
\usepackage{unicode-math}
\setmathfont{latinmodern-math.otf}
\AtBeginDocument{\let\setminus\smallsetminus}
\usepackage{xurl,hyperref}
\hypersetup{colorlinks=true,urlcolor=blue}
\setcounter{secnumdepth}{0}
\setlength{\parindent}{0pt}
\setlength{\parskip}{5pt}
\setlength{\emergencystretch}{3em}
\begin{document}
\section*{Invariant line fields on Julia sets}\label{5300044}
{\small UnsolvedMath ID 5300044 · AMR-052-0044 · Dynamical Systems · AMR Open Problem Lists}
\subsection{Definitions and mathematical statement}
Let $f:\widehat{\mathbb C}\to\widehat{\mathbb C}$
be a rational map of degree at least two.
Its Julia set $J(f)$ is the complement of
the open set on which the iterates form
a normal family. A measurable invariant
line field on $J(f)$ means a measurable
choice of an unoriented real tangent line
$\ell(z)\subset T_z\widehat{\mathbb C}$ on
an invariant measurable subset $E\subset J(f)$
of positive spherical area, such that
\[
                    Df_z(\ell(z))=\ell(f(z))
                           \quad\text{for almost every }z\in E.
\]
Invariance of $E$ is modulo area-zero sets.
Critical points, a finite set, do not
affect the almost-everywhere condition.

A Latt\`es map is a rational map arising
from an affine expanding map $L(w)=aw+b$
of a complex torus $\mathbb C/\Lambda$,
where $a\Lambda\subset\Lambda$ and $|a|>1$,
by a finite holomorphic quotient
$\pi:\mathbb C/\Lambda\to\widehat{\mathbb C}$
satisfying $f\pi=\pi L$.

Does the existence of such a line field
force $f$ to be a Latt\`es map? Equivalently,
prove that every rational map outside
this exceptional class has no measurable
invariant line field supported on a
positive-area part of its Julia set.
The positive-area requirement excludes fields
specified merely on countable or zero-area
invariant sets. No continuity of the
field is assumed, and line fields on
Fatou components are outside the question.
\subsection{Short English statement}
Are rational maps with a measurable invariant line field on a positive-area part of their Julia set necessarily Lattès maps?
\subsection{Sources}
\begin{itemize}
\item[S1] ULAM.AI and the UnsolvedMath Contributors, supplied problems.json, record 5300044 (AMR-052-0044), AMR Open Problem Lists. Catalogue identity and source transcription; CC BY 4.0.
\url{https://huggingface.co/datasets/ulamai/UnsolvedMath}.
\item[S2] Stony Brook - Open Problems in Dynamical Systems, 1992, Lyubich P4. Original source indexed by AMR-052-0044.
\url{https://www.math.stonybrook.edu/open-problems-dynamical-systems}.
\item[S3] Lyubich, Measure and Dimension of Julia Sets, Problem4 (Sullivan), in Problems in Holomorphic Dynamics, IMS92/7, p.27.
\url{https://www.math.stonybrook.edu/preprints/ims92-7.pdf}.
\end{itemize}
\end{document}
